% Theodore von K\'arm\'an, "The engineer grapples with nonlinear problems", Bull. Amer. Math. Soc. 46 (1940), first page, reset in Mills 8A.
\documentclass[11pt]{article}
\usepackage{bulletin1947}
\usepackage{array}
\begin{document}
\setcounter{page}{615}
\articletitle{THE ENGINEER GRAPPLES WITH NONLINEAR PROBLEMS\footnotemark}{theodore von k\'arm\'an}
\footnotetext{The fifteenth Josiah Willard Gibbs Lecture, delivered at Columbus, Ohio, December 27, 1939, under the auspices of the American Mathematical Society with the cooperation of the American Association for the Advancement of Science.}
I do not believe that one could connect justly the name of Gibbs with practical applications of applied mathematics, for his main interest was certainly centered on basic conceptions of mathematical physics. Nevertheless, for example, his beautiful work on graphical methods in thermodynamics is a brilliant example of the presentation of theoretical relations in a form which appeals to the engineer.

This lecture is intended as an effort to improve the convergence between the viewpoints of mathematics and engineering. Thus, I feel it is not inappropriate to dedicate it to the memory of Josiah Willard Gibbs.

Engineering mathematics is generally considered as a collection of mathematical methods adapted for the solution of relatively simple problems. These problems often might require lengthy numerical calculations or graphical constructions, but supposedly can be worked out without the use of advanced methods of mathematical analysis. This description was perhaps correct some decades ago; today a large group of scientific workers is engaged in applying various methods of classical and modern analysis to problems in electrical, civil, mechanical, aeronautical and also chemical engineering. It is not possible to give an exhaustive list of all types of problems which require the applications of advanced analytical methods. In the following table merely some of the most important engineering problems and the mathematical concepts and methods involved in their treatment are indicated:

\begin{center}
\footnotesize\textsc{Associated topics of engineering and mathematics}

\smallskip
\footnotesize
\newcolumntype{P}{>{\setlength\parindent{1.5em}\noindent\hspace{\parindent}}p{\dimexpr.5\textwidth-2\tabcolsep-.2pt\relax}}
\setlength\tabcolsep{4pt}
\begin{tabular}{@{}P|P@{}}
\hline\hline
\multicolumn{1}{@{}c|}{\rule{0pt}{14pt}Mathematical topics} & \multicolumn{1}{c@{}}{Engineering problems}\\[6pt]
\cline{1-1}\cline{2-2}
\rule{0pt}{14pt}Vector algebra, systems of linear equations. Tensors and matrices. Algebraic equations. Ordinary differential equations with given initial conditions. Elementary operational calculus.
& \rule{0pt}{14pt}Engine dynamics, vibration of systems with a finite number of degrees of freedom. Rotating electric machinery.\\[6pt]
\cline{1-1}\cline{2-2}
\rule{0pt}{14pt}Ordinary differential equations and their boundary problems. Eigenvalues and eigenfunctions. Expansion in or
& \rule{0pt}{14pt}Equilibrium, buckling and harmonic vibrations of beams. Critical frequencies and speeds. One-dimensional problems
\end{tabular}
\end{center}
\end{document}
